\documentclass[11pt]{article}
% Shared print layout. The web's provisional design and MathJax are independent.
\usepackage[utf8]{inputenc}
\usepackage[T1]{fontenc}
\usepackage[spanish,es-nodecimaldot]{babel}
\usepackage{lmodern}
\usepackage{amsmath,amssymb}
\usepackage{graphicx}
\usepackage{geometry}
\geometry{a4paper,left=20mm,right=20mm,top=25mm,bottom=22mm,headheight=16pt,headsep=8mm}
\usepackage{fancyhdr,lastpage}
\usepackage{enumitem}
\usepackage{xcolor}
\usepackage{pgfplots}
\pgfplotsset{compat=1.18}
\usepgfplotslibrary{fillbetween}
\usepackage[hidelinks,pdfencoding=auto]{hyperref}
\hypersetup{pdflang=es-ES,pdfauthor={},pdfsubject={Matemáticas Aplicadas a las Ciencias Sociales II}}
\pagestyle{fancy}
\fancyhf{}
\fancyhead[L]{\small Matemáticas Aplicadas a las Ciencias Sociales II}
\fancyhead[R]{\small 2.\textsuperscript{o} Bachillerato}
\fancyfoot[L]{\footnotesize Funciones: variante | Edición revisada 1.0}
\fancyfoot[R]{\footnotesize Página \thepage\ de \pageref{LastPage}}
\renewcommand{\headrulewidth}{0.4pt}
\setlength{\parindent}{0pt}
\setlength{\parskip}{0.5em}
\setlist{itemsep=0.3em,topsep=0.3em}
\newcommand{\apartado}[1]{\subsection*{#1}}
\pgfplotsset{
  revision/.style={
    width=0.94\linewidth,height=7.5cm,
    axis lines=middle,grid=major,grid style={gray!18},
    tick label style={font=\small},label style={font=\small},
    xlabel={$x$},ylabel={$y$},samples=160,
    legend style={at={(0.5,-0.16)},anchor=north,draw=none,font=\small},
    clip=true
  }
}

\hypersetup{pdftitle={Funciones: variante - Soluciones revisadas}}
\begin{document}
\section*{Ejercicio 1: función por tramos}
Solucionario elaborado en esta revisión a partir de F006. El ZIP no incluía soluciones.

El dominio es $[-2,2]$. Las dos ramas son continuas y, en $0$, los límites laterales y $f(0)$ valen 1. Por tanto, es continua en todo su dominio, con continuidad relativa al dominio en los extremos.
\[f'(x)=\begin{cases}2(x+1),&-2<x<0,\\2(x-1),&0<x<2.\end{cases}\]
En $0$, las derivadas laterales son 2 y $-2$: \textbf{no es derivable}. En $-2$ y $2$ no se habla de derivada ordinaria bilateral; las derivadas unilaterales interiores valen $-2$ y $2$, respectivamente.

Decrece en $(-2,-1)$ y $(0,1)$; crece en $(-1,0)$ y $(1,2)$. Mínimos absolutos: $f(-1)=f(1)=0$. Máximo absoluto 1 en $-2,0,2$. El extremo relativo interior máximo es $0$ y los mínimos relativos interiores son $-1$ y $1$. Si se admiten extremos relativos unilaterales, los extremos del dominio también son máximos relativos.

En cada intervalo abierto de rama, $f''=2>0$: curvatura hacia arriba. En $0$ no existe segunda derivada y no cambia el sentido de la curvatura entre las ramas; no es punto de inflexión. \textbf{No es convexa en todo $[-2,2]$}: $f(0)=1>(f(-1)+f(1))/2=0$.
\begin{center}\begin{tikzpicture}\begin{axis}[revision,height=6.5cm,xmin=-2.4,xmax=2.4,ymin=-0.2,ymax=1.4,xtick={-2,-1,0,1,2},ytick={0,1}]
\addplot[thick,domain=-2:0] {(x+1)^2};\addplot[thick,domain=0:2] {(x-1)^2};
\addplot[only marks,mark=*] coordinates {(-2,1) (-1,0) (0,1) (1,0) (2,1)};
\end{axis}\end{tikzpicture}\end{center}
\clearpage
\section*{Ejercicio 2: asíntotas y hueco}
El dominio original excluye $2,3,5$. Solo para $x\neq2$ puede simplificarse
\[f(x)=\frac1{(x-5)(x-3)}.\]
\textbf{En $2$ no hay asíntota vertical}: el límite es $1/3$, finito. Existe un hueco en $(2,1/3)$; simplificar no reincorpora ese punto al dominio.
\begin{center}\begin{tabular}{c|cc}
Valor & Límite por la izquierda & Límite por la derecha\\\hline
$3$ & $+\infty$ & $-\infty$\\
$5$ & $-\infty$ & $+\infty$
\end{tabular}\end{center}
Asíntotas verticales: $x=3$ y $x=5$. Como el denominador simplificado tiene grado 2, $\lim_{x\to\pm\infty}f(x)=0$: asíntota horizontal $y=0$. No hay oblicua.
\begin{center}\begin{tikzpicture}\begin{axis}[revision,height=7cm,xmin=-1,xmax=8,ymin=-3,ymax=3,xtick={0,2,3,4,5,6,8}]
\addplot[thick,domain=-1:1.99] {1/((x-5)*(x-3))};\addplot[thick,domain=2.01:2.97] {1/((x-5)*(x-3))};
\addplot[thick,domain=3.03:4.97] {1/((x-5)*(x-3))};\addplot[thick,domain=5.03:8] {1/((x-5)*(x-3))};
\addplot[dashed] coordinates {(3,-3) (3,3)};\addplot[dashed] coordinates {(5,-3) (5,3)};
\addplot[only marks,mark=o,mark size=3pt,mark options={fill=white}] coordinates {(2,0.333333333)};
\end{axis}\end{tikzpicture}\end{center}
Las ramas se dibujan separadamente, sin evaluar los polos. El círculo vacío señala el hueco, no una asíntota.
\clearpage
\section*{Ejercicio 3: signo de la derivada}
Factorizamos exactamente:
\[c'(t)=\frac3{100}(t-10)(t-20).\]
Los puntos críticos interiores son $t=10$ y $t=20$.
\begin{center}\begin{tabular}{c|ccc}
Intervalo & $(0,10)$ & $(10,20)$ & $(20,24)$\\\hline
Signo de $c'$ & $+$ & $-$ & $+$\\
Comportamiento de $c$ & Creciente & Decreciente & Creciente
\end{tabular}\end{center}
En $10$ cambia de creciente a decreciente: \textbf{máximo relativo}. En $20$ cambia de decreciente a creciente: \textbf{mínimo relativo}.

Los extremos $0$ y $24$ no son puntos críticos interiores. No se inventan cotizaciones concretas: conocer $c'$ determina $c$ salvo una constante, pero no su nivel de precios.
\begin{center}\begin{tikzpicture}\begin{axis}[revision,height=7cm,xmin=0,xmax=24,ymin=-1.2,ymax=6.8,xtick={0,10,15,20,24},ylabel={$c'(t)$},xlabel={Horas ($t$)}]
\addplot[thick,domain=0:24] {0.03*(x-10)*(x-20)};
\addplot[only marks,mark=*] coordinates {(10,0) (20,0)};
\end{axis}\end{tikzpicture}\end{center}
La figura representa la \textbf{derivada}, no el precio $c$. La variante bursátil que aparece comentada en ambos ZIP se conserva una sola vez en el anexo no evaluable del solucionario principal.
\end{document}
